\documentclass[aspectratio=169,xcolor=dvipsnames]{beamer}

\usepackage{amsthm}
\usepackage{amsmath}
\usepackage{amssymb}
\usepackage{tikz}
\usepackage{pgfplots}
\usepackage{xfp}
\pgfplotsset{compat=1.18}
\usepackage[authoryear, round]{natbib}
\usepackage{hyperref}
\usetikzlibrary{arrows.meta,positioning,decorations.pathreplacing}
\hypersetup{
    colorlinks=true,
    linkcolor=blue,
    citecolor=blue,
    urlcolor=blue
}

\definecolor{darkgreen}{rgb}{0,0.5,0}
\definecolor{darkblue}{rgb}{0,0,0.55}

\newtheorem{proposition}{Proposition}
\setbeamertemplate{footline}[frame number]
\usetheme{Frankfurt}
\usefonttheme{serif}

% Title info
\title{ECON8037: Financial Economics}
\subtitle{\normalsize Week 8 Lecture - The Economics of Cryptocurrencies II: Optimal Design, Welfare, and Payments (Chiu \& Koeppl, 2017)}
\author{Simon Mishricky}
\institute{Research School of Economics, Australian National University}
\date{\today}

\begin{document}
\section{Recap}
\frame{\titlepage}

\begin{frame}{Where we left off}
    Lecture I built a Lagos--Wright model with a cryptocurrency:
    \vspace{2mm}
    \begin{itemize}
        \item Mining is competitive and dissipates all rewards: $MQ = \beta R$.
        \item Buyers can attack via secret mining. A contract $(d,N)$ is
              \emph{double-spending proof} iff
              \[
                  d \;<\; R(N+1)N.
              \]
        \item Settlement cannot be both immediate and final.
        \item Mining reward is endogenous:
              \[
                  R = \frac{Z(\mu-1) + D\tau}{\bar N+1}.
              \]
    \end{itemize}
    \vspace{2mm}
    \textbf{Today.} Optimal $(\mu,\tau)$, calibration to Bitcoin, welfare, comparison with
    Fedwire and US debit cards.
\end{frame}

\begin{frame}{Roadmap for Lecture II}
    \begin{enumerate}
        \item Optimal cryptocurrency design (the planner's Ramsey problem).
        \item Calibrating the model to Bitcoin.
        \item Welfare cost: Bitcoin vs.\ cash vs.\ optimal cryptocurrency.
        \item Best use cases: Fedwire (large-value) vs.\ US debit cards (retail).
        \item Scalability and policy implications.
    \end{enumerate}
\end{frame}

\section{Optimal design}

\begin{frame}{Two design instruments}
    The cryptocurrency designer chooses $(\mu,\tau)$:
    \vspace{2mm}
    \begin{itemize}
        \item $\mu \ge 1$ -- gross money growth rate (\emph{seigniorage}).
        \item $\tau \ge 0$ -- transaction fee (fraction of $d$ paid to miners).
    \end{itemize}
    \vspace{3mm}
    Both finance mining rewards (which dissipate as mining cost):
    \[
        \frac{\beta}{\mu} R(\bar N+1) \;=\; B(\mu - 1 + \sigma\tau)\!\int \frac{x(\varepsilon)}{1-\tau}\,dF_\varepsilon(\varepsilon).
    \]
    \vspace{1mm}
    \textbf{Trade-off.} More mining $\Rightarrow$ tighter DS-proof constraint $\Rightarrow$ larger trade.
    But mining is a deadweight loss and inflation/fees distort consumption.
\end{frame}

\begin{frame}{Social welfare}
    Welfare = average utility per period:
    \[
        \mathcal{W} \;=\; B\sigma\!\int_0^\infty \bigl[\delta^{N(\varepsilon)}\varepsilon u(x(\varepsilon)) - x(\varepsilon)\bigr] dF_\varepsilon(\varepsilon) - C(\bar N+1). \tag{37}
    \]
    \vspace{1mm}
    By Lemma 1 (mining dissipation), the mining cost equals $\beta R(\bar N+1)/\mu$,
    the present value of total rewards.
\end{frame}

\begin{frame}{The planner's problem}
    \[
        \max_{\mu,\tau} \;\; B\!\int \!\sigma\,\delta^{N(\varepsilon)}\varepsilon u[x(\varepsilon)] - x(\varepsilon)\,dF_\varepsilon(\varepsilon) - \tfrac{\beta}{\mu}R(\bar N+1) \tag{38}
    \]
    subject to
    \begin{itemize}
        \item $(d(\varepsilon),x(\varepsilon),N(\varepsilon))$ is the buyer's optimal contract given $(\mu,\tau)$;
        \item DS constraint: $d(\varepsilon) = \frac{\mu}{\beta}\frac{x(\varepsilon)}{1-\tau} \le R(N(\varepsilon)+1)N(\varepsilon)$;
        \item Planner's budget: $\tfrac{\beta}{\mu}R(\bar N+1) = B(\mu - 1 + \sigma\tau)\!\int\!\tfrac{x(\varepsilon)}{1-\tau}\,dF_\varepsilon(\varepsilon)$.
    \end{itemize}
    \vspace{2mm}
    This is a \emph{Ramsey problem}: planner internalises that agents respond to $(\mu,\tau)$.
\end{frame}

\begin{frame}{The optimal design result}
    \begin{proposition}[Optimal design]
        The optimal reward structure sets transaction fees to zero and relies only on
        seigniorage:
        \[
            \boxed{\;\tau = 0 \quad\text{and}\quad \mu > 1.\;}
        \]
    \end{proposition}
    \vspace{3mm}
    \textbf{Why?} Both $\mu$ and $\tau$ distort consumption, but their incidence differs:
    \begin{itemize}
        \item Inflation tax is paid by \emph{all} balance holders.
        \item Transaction fees are paid only by buyers who \emph{actually trade} -- precisely
              the agents with the highest valuation of balances.
    \end{itemize}
    Levying revenue through inflation smooths the distortion across all buyers.
\end{frame}

\begin{frame}{Intuition for the trade-off}
    \begin{itemize}
        \item Without rewards: no mining, no security, no equilibrium.
        \item With too many rewards: mining is excessive; large deadweight loss.
        \item Optimum balances the marginal benefit of relaxing the DS constraint
              against the marginal cost of mining.
    \end{itemize}
    \vspace{3mm}
    The result aligns with the broader monetary literature: \emph{ex ante} taxation (inflation
    on all balances) dominates \emph{ex post} user fees when the latter target the most
    elastic margin.
\end{frame}

\section{Calibration}

\begin{frame}{Calibrating the model to Bitcoin}
    Utility $u(x) = \log(x+b) - \log b$ with $b\approx 0$.
    A period = 1 day, a trading session = 10 minutes (the average block time).
    \vspace{2mm}
    \begin{itemize}
        \item $\beta = 0.9999$ ($\approx 0.97$ annual);
        \item 2015 Bitcoin supply: $14{,}342{,}503$ BTC;
        \item daily money growth: $\mu = (1 + 25/14342503)^{6\times 24} = 1.00025$, $\approx 9.6\%$ annual;
        \item $\sigma = 0.0178$ (fraction of BTC spent per day);
        \item $\tau = 0.0088\%$ (from data on fees);
        \item $B = 6{,}873{,}428$ (max number of average-sized transactions).
    \end{itemize}
\end{frame}

\begin{frame}{Benchmark parameters}
    \centering
    \footnotesize
    \begin{tabular}{lll}
        \hline
        Parameter & Value & Target \\
        \hline
        $\beta$ & 0.9999166 & period = 1 day \\
        $\delta$ & 0.9999994 & $\delta = \beta^{1/(1+\bar N)}$ \\
        $\mu$ & 1.00025 & Bitcoin money growth \\
        $\tau$ & 0.000088 & Bitcoin transaction fee \\
        $B$ & 6{,}873{,}428 & max no.\ of average-sized transactions \\
        $\sigma$ & 0.0178 & velocity per block \\
        $\alpha$ & 1 & normalisation \\
        \hline
    \end{tabular}
    \vspace{4mm}

    Preference shock distribution $F(\varepsilon)$ disciplined by the empirical distribution
    of transaction sizes in \citet{ron2013}.
\end{frame}

\begin{frame}{Effects of money growth}
    Numerical experiment around the calibrated $\mu$:
    \vspace{2mm}
    \begin{itemize}
        \item Higher $\mu$ $\Rightarrow$ higher $R$ $\Rightarrow$ higher mining effort,
              shorter average confirmation lag $N$.
        \item Aggregate trade and utility \emph{rise} (DS constraint relaxes, lags shrink).
        \item Mining cost \emph{rises} linearly with $\mu$.
        \item Net welfare: \textbf{hump-shaped} in $\mu$.
    \end{itemize}
    \vspace{3mm}
    Bitcoin's calibrated $\mu$ is on the \emph{high} side: more inflation than is optimal.
\end{frame}

\section{Welfare}

\begin{frame}{Welfare measures}
    Two consumption-equivalent measures:
    \vspace{2mm}
    \begin{itemize}
        \item \textbf{Welfare cost.} Fraction of consumption a person would give up to use
              cash under the Friedman rule (zero inflation) instead of cryptocurrency.
        \item \textbf{Equivalent inflation tax.} Cash inflation rate that makes a person
              indifferent between traditional cash and the cryptocurrency.
    \end{itemize}
\end{frame}

\begin{frame}{Bitcoin is very inefficient}
    \centering
    \small
    \begin{tabular}{lcc}
        \hline
        & Bitcoin (benchmark) & Bitcoin (optimal) \\
        \hline
        $\mu - 1$ & 9.5\% & 0.17\% \\
        $\tau$ & 0.0088\% & 0\% \\
        welfare loss & \textbf{1.41\%} & \textbf{0.08\%} \\
        mining cost (per year) & \$359.98 mn & \$6.90 mn \\
        equivalent inflation tax & 230.44\% & 27.51\% \\
        \hline
    \end{tabular}
    \vspace{4mm}
    \begin{itemize}
        \item Bitcoin's welfare cost: $\approx 1.4\%$ of consumption $\approx$ 500$\times$ a low-inflation
              cash system (whose welfare cost is $\approx 0.003\%$).
        \item Equivalent inflation tax of \textbf{230\%}: people would prefer cash even at
              hyperinflationary rates.
    \end{itemize}
\end{frame}

\begin{frame}{Why is Bitcoin so inefficient?}
    The driver is the \emph{mining cost} -- about \$360 million per year at the calibrated parameters.
    \vspace{3mm}
    \begin{itemize}
        \item Bitcoin's reward structure (high $\mu$ + positive $\tau$) overcompensates miners.
        \item Optimal design (Proposition: $\tau=0$, lower $\mu$) reduces mining to
              \$6.9 million per year.
        \item Welfare loss falls from 1.41\% to \textbf{0.08\%} -- comparable to a cash system with
              moderate inflation.
    \end{itemize}
    \vspace{3mm}
    The lesson: Bitcoin's inefficiency is largely a \emph{design} problem, not an
    intrinsic feature of cryptocurrencies.
\end{frame}

\section{Best use cases}

\begin{frame}{Retail vs.\ large value payments}
    Recalibrate the model to two real-world payment systems (2014 US):
    \vspace{2mm}
    \begin{itemize}
        \item \textbf{US Debit cards} (retail): avg transaction \$38.29; annual volume 59{,}539 mn.
        \item \textbf{Fedwire} (large value): avg transaction \$6.55 million; annual volume 135 mn.
    \end{itemize}
    \vspace{3mm}
    Period = 30 mins; block length = 1 min. Choose $\sigma$ and $\varepsilon$ to match each system.
\end{frame}

\begin{frame}{Welfare comparison}
    \centering
    \small
    \begin{tabular}{lcc}
        \hline
        & Retail (Debit) & Large value (Fedwire) \\
        \hline
        avg transaction size & \$38.29 & \$6{,}552{,}236 \\
        annual volume & 59{,}539 mn & 135 mn \\
        optimal $\mu - 1$ & 0.038\% & 0.53\% \\
        optimal $\tau$ & 0\% & 0\% \\
        confirmation lag & 2 mins & 12 mins \\
        \textbf{welfare loss} & \textbf{0.00052\%} & \textbf{0.0060\%} \\
        mining cost (per year) & \$4.33 mn & \$22.10 bn \\
        equivalent fee (per transfer) & \$0.0002 & \$392.56 \\
        \hline
    \end{tabular}
    \vspace{2mm}
    \begin{itemize}
        \item Optimal $\mu$ is much lower than Bitcoin's for both systems.
        \item Retail welfare loss is two orders of magnitude smaller than large-value.
    \end{itemize}
\end{frame}

\begin{frame}{Why is retail more efficient?}
    \begin{itemize}
        \item Double-spending incentives \emph{scale with transaction size} $d$, while mining cost
              is a public good financed by aggregate volume.
        \item Small individual transactions $\Rightarrow$ DS constraint is easy to satisfy with little mining.
        \item Large individual transactions $\Rightarrow$ require lots of mining and long lags.
    \end{itemize}
    \vspace{3mm}
    Equivalent fee comparison:
    \begin{itemize}
        \item Cryptocurrency retail: 0.02 cents/transfer vs.\ debit card interchange $\approx$ 23 cents.
        \item Cryptocurrency large value: \$392 vs.\ Fedwire service fee $\approx$ 82 cents.
    \end{itemize}
    \vspace{3mm}
    Cryptocurrencies can plausibly compete with retail systems; for large-value
    transactions, centralised systems remain far more efficient.
\end{frame}

\section{Policy and scaling}

\begin{frame}{Scalability is the binding constraint}
    For cryptocurrencies to compete with debit cards they must process \emph{enormous} transaction
    volumes.
    \vspace{2mm}
    \begin{itemize}
        \item Bitcoin: $\approx 7$ transactions per second (TPS) cap.
        \item Visa: $\approx 1{,}700$ TPS sustained, peaks $> 24{,}000$.
        \item Block-size limits, network latency, propagation delays all bind.
    \end{itemize}
    \vspace{3mm}
    The model takes scaling as exogenous, but the welfare results say:
    \emph{if technological scalability is solved, optimally designed cryptocurrencies
    can challenge retail payment networks.}
\end{frame}

\begin{frame}{Caveats}
    The comparison is suggestive, not definitive:
    \vspace{2mm}
    \begin{itemize}
        \item Ignores private costs: data storage, wallets, exchanges, software.
        \item Mining cost is a deadweight loss; some of Visa/Fedwire's fees are
              provider \emph{profits}, not operating cost.
        \item Ignores network/scaling constraints (assumed solved).
        \item Ignores volatility of the cryptocurrency's value.
    \end{itemize}
\end{frame}

\begin{frame}{Policy implications}
    \begin{itemize}
        \item \textbf{Central bank digital currency (CBDC).} The model puts bounds on
              the cost of issuing digital currency centrally vs.\ via PoW.
        \item \textbf{Reward design.} Existing cryptocurrencies should consider eliminating
              transaction fees and using only seigniorage.
        \item \textbf{Reform Bitcoin?} Lower money growth and remove fees $\Rightarrow$ welfare cost
              drops by an order of magnitude.
        \item \textbf{Scaling debate.} On-chain scaling is the binding technological constraint
              if cryptocurrencies are to displace cards.
    \end{itemize}
\end{frame}

\section{Wrap-up}

\begin{frame}{What we've learned}
    \begin{enumerate}
        \item \textbf{Optimal design.} Finance mining via seigniorage only ($\tau=0$, $\mu>1$).
              Inflation spreads the distortion across all balance holders.
        \vspace{2mm}
        \item \textbf{Bitcoin is poorly designed.} Welfare cost 1.4\% of consumption,
              mostly wasted on mining.
        \vspace{2mm}
        \item \textbf{Best use case is retail.} Cryptocurrencies scale efficiently in
              volume but not in transaction size.
        \vspace{2mm}
        \item \textbf{Scaling is the real bottleneck.} If technological limits can be
              overcome, well-designed cryptocurrencies could compete with debit cards.
    \end{enumerate}
\end{frame}

\begin{frame}{Open questions and extensions}
    \begin{itemize}
        \item Endogenous scalability: model block size, latency, network propagation.
        \item Alternative consensus protocols (proof-of-stake, BFT).
        \item Volatility and price discovery for the cryptocurrency.
        \item Interaction with central bank policy (CBDCs, negative rates).
        \item Coordinated double-spending attacks (footnote 21: $d n_b < R(N+1)N$).
        \item Smart contracts and programmable money.
    \end{itemize}
\end{frame}

{
\setbeamercolor{background canvas}{bg=black}
\begin{frame}
    \centering
    \vspace{2mm}
    \Huge \textcolor{white}{\textbf{Thank you}}
\end{frame}
}

\begin{frame}[allowframebreaks]{References}
    \small
    \begin{thebibliography}{99}
        \bibitem[Chiu and Koeppl(2017)]{chiu2017}
            Chiu, J.\ and T.\ V.\ Koeppl (2017). The Economics of Cryptocurrencies --
            Bitcoin and Beyond. \textit{Bank of Canada Staff Working Paper / Queen's University}.
        \bibitem[Lagos and Wright(2005)]{lagos2005}
            Lagos, R.\ and R.\ Wright (2005). A Unified Framework for Monetary Theory
            and Policy Analysis. \textit{Journal of Political Economy} 113(3), 463--484.
        \bibitem[Ron and Shamir(2013)]{ron2013}
            Ron, D.\ and A.\ Shamir (2013). Quantitative Analysis of the Full Bitcoin Transaction Graph.
            \textit{Financial Cryptography}.
    \end{thebibliography}
\end{frame}

\end{document}
